%! Setting Up Determinants --- the 2 by 2 Determinant, Area, and Orientation — Unit 5, Lesson 5.0 — Group Activity
\documentclass[10pt]{article}
\usepackage{linalg-article}
\usepackage{linalg-boxes}
\newcommand{\TallMath}[1]{$\displaystyle #1\rule[-1.4em]{0pt}{3.2em}$}
\newcommand{\vv}{\mathbf{v}}
\newcommand{\ww}{\mathbf{w}}
\newcommand{\avec}{\mathbf{a}}
\newcommand{\bb}{\mathbf{b}}
\newcommand{\zero}{\mathbf{0}}

\begin{document}

\pageheader{Unit 5, Lesson 5.0}{Group Activity}
\namepartnerperiod

\vspace{0.10in}

\begin{headlinebox}{blush}
\small\textbf{Area, Orientation, and the Vanishing Box.} One number, $\det\begin{bsmallmatrix} a & b\\ c & d \end{bsmallmatrix}=ad-bc$,
reads off the \textbf{area} the columns span ($|\det|$), the \textbf{orientation} (its sign), and
whether the matrix is \textbf{invertible} ($\det\ne0$). Work with your partner, show each
determinant, and \emph{justify} every claim.
\end{headlinebox}

\vspace{0.10in}

\begin{tcolorbox}[colback=white, colframe=black!40, title=\textbf{Tier R --- Remediate},
    fonttitle=\bfseries, arc=1.5mm, left=3mm, right=3mm, top=2mm, bottom=2mm, breakable]
Compute and read off area. Show your work.
\begin{enumerate}[leftmargin=1.7em, itemsep=8pt]
  \item Compute the determinant
        $\det\begin{bsmallmatrix} 4 & 1\\ 2 & 3 \end{bsmallmatrix} = \blank{2.4cm}$.
  \item Find the \textbf{area} of the parallelogram with columns
        $\begin{bsmallmatrix} 4\\0 \end{bsmallmatrix}$ and $\begin{bsmallmatrix} 1\\3 \end{bsmallmatrix}$:
        \quad area $=|\det| = \blank{2.4cm}$.
  \item Is the box \emph{flat}? Compute $\det\begin{bsmallmatrix} 2 & 3\\ 4 & 6 \end{bsmallmatrix}$ and
        say whether the columns are parallel. \; \blank{3.0cm}
\end{enumerate}
\end{tcolorbox}

\begin{tcolorbox}[colback=white, colframe=black!40, title=\textbf{Tier A --- Approaching Proficiency},
    fonttitle=\bfseries, arc=1.5mm, left=3mm, right=3mm, top=2mm, bottom=2mm, breakable]
Now use the \emph{sign} and connect to \emph{invertibility}.
\begin{enumerate}[leftmargin=1.7em, itemsep=9pt]
  \item \textbf{Orientation.} Compute $\det\begin{bsmallmatrix} 1 & 4\\ 2 & 3 \end{bsmallmatrix}$. Is the
        orientation \emph{kept} or \emph{flipped}? What is the area? \writelines{2}
  \item \textbf{The swap rule.} Swap the two columns to get $\begin{bsmallmatrix} 4 & 1\\ 3 & 2 \end{bsmallmatrix}$
        and compute its determinant. How does it compare to item~1, and what changed about the box? \writelines{2}
  \item \textbf{Solve for a flat box.} Find $x$ so that $A=\begin{bsmallmatrix} 2 & 6\\ 3 & x \end{bsmallmatrix}$
        has $\det=0$. Is $A$ invertible for that $x$? Explain using the columns. \writelines{2}
\end{enumerate}
\end{tcolorbox}

\begin{tcolorbox}[colback=white, colframe=black!40, title=\textbf{Tier E --- Extension},
    fonttitle=\bfseries, arc=1.5mm, left=3mm, right=3mm, top=2mm, bottom=2mm, breakable]
\textbf{A first look at transformations.} Applying a matrix $A$ to the \textbf{unit square} (corners
$(0,0),(1,0),(1,1),(0,1)$, area $1$) sends it to the parallelogram spanned by $A$'s columns. Its area
is $|\det A|$ --- so $|\det A|$ is the factor \emph{every} area gets multiplied by (Lesson~5.3).
\begin{enumerate}[leftmargin=1.7em, itemsep=9pt]
  \item Let $A=\begin{bsmallmatrix} 2 & 1\\ 0 & 3 \end{bsmallmatrix}$. The image of the unit square is
        spanned by its columns $\begin{bsmallmatrix} 2\\0 \end{bsmallmatrix}$ and
        $\begin{bsmallmatrix} 1\\3 \end{bsmallmatrix}$. Find the image's area. By what factor did the
        area grow? \writelines{2}
  \item Now let $B=\begin{bsmallmatrix} 1 & 2\\ 2 & 1 \end{bsmallmatrix}$. Compute $\det B$. What is the
        area factor, and what does the \emph{sign} say happened to the square? \writelines{2}
  \item A matrix has $\det=0$. What happens to the unit square when you apply it? Why does that mean
        the transformation cannot be \emph{undone}? \writelines{2}
\end{enumerate}
\end{tcolorbox}

\end{document}
